\documentclass[10pt,a4paper]{amsart}
\usepackage[margin=19mm]{geometry}
\usepackage{amsmath,amssymb,mathtools}
\usepackage[scr=rsfs,cal=euler]{mathalfa}
\usepackage{xfrac}
\usepackage[unicode,hidelinks,bookmarks=false]{hyperref}
\newcommand{\ie}{i.e.\ }
\newcommand{\cf}{cf.\ }
\newcommand{\resp}{respectively\ }
\newcommand{\Subsection}{\subsection}
\providecommand{\nobreakdash}{\nobreak\hbox{-}}
\setlength{\emergencystretch}{2em}
\multlinegap0pt
\usepackage{siunitx}
\usepackage{siunitx}
\usepackage{siunitx}
\usepackage{siunitx}
\usepackage{siunitx}
\usepackage{csquotes}
\usepackage{tikz}
\newtheorem{lem}{Lemma}
\newcommand{\st}{\text{s.t.}}
\title{Mixed-Integer Linear Optimization for Semi-Supervised Optimal Classification Trees}
\author{Jan Pablo Burgard, Maria Eduarda Pinheiro, Martin Schmidt}
\date{Source: arXiv:2401.09848v3 (CC BY 4.0)}
\begin{document}
\maketitle

\noindent\emph{Source and licence.} Mixed-Integer Linear Optimization for Semi-Supervised Optimal Classification Trees. Version arXiv:2401.09848v3 (CC BY 4.0), distributed by arXiv under the Creative Commons Attribution 4.0 International licence, http://creativecommons.org/licenses/by/4.0/, as recorded in the arXiv licence metadata for this exact version and shown on the abstract page https://arxiv.org/abs/2401.09848v3. Full licence notice: \texttt{licenses/arxiv-2401.09848-CC-BY-4.0.txt}.\par\smallskip

\noindent\textit{Editorial scope.} Counted unit: the article's lem lemma2 (source0.tex lines 703--749). The article's complete original proof of it is delivered with it (source0.tex lines 750--775, one \texttt{proof} environment of 26 lines). No mathematics is written, repaired or re-derived: the statement and the proof are the article's own lines, in the article's own notation, and no retained line is substituted or re-flowed.  No further source-local context is needed: every cross-reference of every retained line resolves inside the two delivered spans.  Nothing was omitted from any retained span, so the package carries no omission record.  No retained line contains a citation, so the wrapper carries no bibliography and no bibliography entry.  No retained line contains a figure environment, an image-inclusion command or drawing code, so no image is delivered and the package declares no asset dependency.  Editorial formatting only, in addition: an AMS \texttt{amsart} preamble that restates the article's own declarations and macros used by the retained lines, namely the environments \texttt{lem} on separate counters, together with the standard packages those lines need.  The retained environments print in this document as Lemma 1 in order of appearance, whereas the article prints the same items with the numbering of its own full document.  The complete native source of the article as distributed in the arXiv e-print for this exact version is bundled unchanged as \texttt{sources/153143/source0.tex}.
\begin{lem}
  \label{lemma2}
  Consider a set of continuous functions $f_k : \mathbb{R}^p \to
  \mathbb{R}$, $k \in [1,d]$, for some $d\in \mathbb{N}$, and let $\Omega
  \subseteq \mathbb{R}^p$ be given.
  Suppose further that there exist values $u_k>0$ such that $$0
  \leq f_k(x)\leq u_k$$ holds for all $x \in  \Omega$ and $k \in
  [1,d]$.
  Then, $x^*\in \mathbb{R}^p$ is a solution to the problem
  \begin{subequations} \label{problemmain}
    \begin{align}
      \min_{x} \quad
      & \min_{k \in [1,d]} f_k(x)
        \label{eq:problemmain1}\\
      \st \quad
      & x \in \Omega
        \label{eq:problemmain2}
    \end{align}
  \end{subequations}
  if and only if there exist $\alpha^*, \beta^* \in \mathbb{R}^d$ such
  that $(x^*, \alpha^*, \beta^*)$ is a solution to the problem
  \begin{subequations}
    \label{problembeta}
    \begin{align}
      \min_{ x,\alpha,\beta }
      & \quad \sum_{k=1}^d \beta_k
        \label{eq:beta1} \\
      \st
      & \quad x \in  \Omega,
        \label{eq:beta2}  \\
      & \quad  \sum_{k =1}^d \alpha_k = 1,
        \label{eq:beta3} \\
      &  \quad    \alpha_k \in \{0,1\}, \quad k \in[1,d],
        \label{eq:beta4}     \\
      &  \quad    \beta_k\leq u_k\alpha_k, \quad k \in[1,d],
        \label{eq:beta5}   \\
      &  \quad    \beta_k\leq f_k(x), \quad k \in[1,d], \quad x \in
        \Omega ,
        \label{eq:beta6}   \\
      &  \quad   \beta_k\geq f_k(x) - u_k(1-\alpha_k), \quad k
        \in[1,d], \quad x \in  \Omega,
        \label{eq:beta7}   \\
      &  \quad    \beta_k\geq 0, \quad k \in[1,d].
        \label{eq:beta8}
    \end{align}
  \end{subequations}
\end{lem}

\begin{proof}
  By introducing the binary variables~$\alpha_k$, we obtain that $x^*$
  is a solution to Problem~\eqref{problemmain} if and only if $(x^*,
  \alpha^*)$ is a solution of the problem
  \begin{align*}
    \min_{ x,  \alpha  } \quad
    & \sum_{k =1}^d  \alpha_k f_k(x)
    \\
    \st \quad
    &  \eqref{eq:beta2}\text{--}\eqref{eq:beta4}.
  \end{align*}

  Besides that, for any $(x^*, \alpha^*)$ solution of the problem
  above, if $\alpha_k^* = 0$ holds for some $k\in [1,d]$, Constraints
  \eqref{eq:beta5} and \eqref{eq:beta8} enforce that $\beta^*_k = 0 =
  \alpha_k f_k(x^*)$. On the other hand, if $\alpha_k^*=1$ holds, by
  Constraints~\eqref{eq:beta6} and~\eqref{eq:beta7}, we obtain
  $\beta_k^* = \alpha^*_k f_k(x^*)$.

  Hence, $x^*$ is a solution to Problem~\eqref{problemmain} if and
  only if $(x^*, \alpha^*,\beta^*)$ exists such that
  \begin{equation*}
    \beta^*_k = \alpha^*_k f_k(x^*), \quad k \in [1,d],
  \end{equation*}
  is a solution to Problem~\eqref{problembeta}.
\end{proof}

\end{document}
